% TT30041C
% Verletzungen mit chemischen Substanzen
% 14.0
% 2013-11-30
\label{m:chemische-verletzungen}
\begin{enumerate}
  \item Auf Selbsschutz achten!
  \begin{itemize}
    \item Patienten ggfs. retten, aber nur wenn es der
    Eigenschutz zulässt.
  \end{itemize}
  \item Gefahr erkennen: \EE{Welcher Stoff?}  Expertenrat einholen!
  \item Wenn erforderlich: Vorgehen analog zu Abschnitt
  \enquote{Gefahrengutunfall}, GAS-Regel
  (\prettyref{topic:gefahrengutunfall})
  \item Weitere Opfer ausschließen (lassen)
  \item Sicherung der Vitalfunktionen, Beurteilung ob der
  Patient vital bedroht ist.
  \item \TxBeiVit \TxMassVitMKBes
  \begin{itemize}
    \item Lagerung: situationsgerecht.
  \end{itemize}
  \item Entfernung kontaminierter Kleidung
  \item Spülung der betroffenen Hautpartien mit Wasser  (richtig abrinnen
  lassen)
  \item Auge: Spülung mit Wasser / Ringerlösung (nach außen abrinnen
  lassen!),  Kontaktlinsen entfernen
  \item Transportentscheidung: Ab einer Ausdehnung von
  15\PC* (Erwachsener) bzw. 10\PC* (Kind) Spezialabteilung
  für Verbrennungsverletzungen. Bei Augenverletzungen Abt.
  für Augenheilkunde.
\end{enumerate}